\documentclass[12pt]{article}

% Packages
\usepackage[a4paper, margin=1in]{geometry}
\usepackage{setspace}
\usepackage{amsmath, amssymb}
\usepackage{hyperref}
\usepackage{titlesec}
\usepackage{enumitem}

% Formatting
\onehalfspacing
\setlist[itemize]{noitemsep}
\setlist[enumerate]{noitemsep}

\titleformat{\section}{\large\bfseries}{\thesection.}{0.5em}{}
\titleformat{\subsection}{\normalsize\bfseries}{\thesubsection}{0.5em}{}

\begin{document}

\begin{center}
{\Large \textbf{CSE400 -- Fundamentals of Probability in Computing}}\\[0.2cm]
{\large Lecture 11: Transformation of Random Variables}\\[0.3cm]

\textbf{Instructor:} Dhaval Patel, PhD\\
\textbf{Date:} February 10, 2026
\end{center}

\hrule
\vspace{0.4cm}

\section{Transformation of Random Variables}

\subsection{Objective}

The goal of transformation techniques is to determine the distribution of a new random variable obtained by applying a function to an existing random variable.

Let

\[
Y = g(X)
\]

where

\begin{itemize}
\item $X$ is a random variable with known distribution
\item $g(\cdot)$ is a deterministic function
\item $Y$ is the transformed random variable
\end{itemize}

The main problem is:

\textbf{Given the distribution of $X$, find the distribution of $Y$.}

\subsection{CDF Method for Transformation}

The most general technique for transformation is the \textbf{CDF method}.

\subsubsection*{Step 1 --- Define the CDF of the new variable}

\[
F_Y(y) = P(Y \le y)
\]

Since $Y = g(X)$,

\[
F_Y(y) = P(g(X) \le y)
\]

Thus, the problem reduces to evaluating this probability in terms of $X$.

\subsection{Case 1 --- Monotonic Increasing Transformation}

Assume $g(x)$ is \textbf{monotonically increasing}, meaning

\[
g(x_1) < g(x_2) \quad \text{if} \quad x_1 < x_2
\]

\subsubsection*{Step 1: Express the event}

\[
F_Y(y) = P(g(X) \le y)
\]

Because $g$ is increasing,

\[
g(X) \le y \iff X \le g^{-1}(y)
\]

Thus,

\[
F_Y(y) = P(X \le g^{-1}(y))
\]

\subsubsection*{Step 2: Use the CDF of $X$}

\[
F_Y(y) = F_X(g^{-1}(y))
\]

\subsubsection*{Step 3: Obtain the PDF}

If $X$ has PDF $f_X(x)$,

\[
f_Y(y) = \frac{d}{dy}F_Y(y)
\]

Substituting,

\[
F_Y(y) = F_X(g^{-1}(y))
\]

Using the chain rule,

\[
f_Y(y) = f_X(g^{-1}(y)) 
\left|\frac{d}{dy} g^{-1}(y)\right|
\]

\subsection{Case 2 --- Monotonic Decreasing Transformation}

If $g(x)$ is \textbf{monotonically decreasing}, the inequality reverses.

Start with

\[
F_Y(y) = P(Y \le y)
\]

Substitute $Y = g(X)$:

\[
F_Y(y) = P(g(X) \le y)
\]

Since $g$ is decreasing,

\[
g(X) \le y \iff X \ge g^{-1}(y)
\]

Thus

\[
F_Y(y) = P(X \ge g^{-1}(y))
\]

Using the CDF definition

\[
P(X \ge a) = 1 - F_X(a)
\]

Therefore

\[
F_Y(y) = 1 - F_X(g^{-1}(y))
\]

Differentiating,

\[
f_Y(y) = \frac{d}{dy}F_Y(y)
\]

which leads to

\[
f_Y(y) = f_X(g^{-1}(y))
\left|\frac{d}{dy}g^{-1}(y)\right|
\]

The absolute value ensures the density remains non-negative.

\section{Function of Two Random Variables}

Consider transformations involving two random variables.

\[
Z = g(X,Y)
\]

where $X$ and $Y$ are random variables with joint distribution

\[
f_{X,Y}(x,y)
\]

The objective is to determine the distribution of the new variable $Z$.

\subsection{Joint Distribution}

If $X$ and $Y$ are continuous random variables, they have a joint probability density function $f_{X,Y}(x,y)$.

The probability of a region $A$ is

\[
P((X,Y)\in A) =
\int\int_A f_{X,Y}(x,y)\,dx\,dy
\]

\section{Illustrative Example}

Consider

\[
Z = X + Y
\]

The goal is to determine the distribution of $Z$.

\subsection{Step 1 --- Define the CDF}

\[
F_Z(z) = P(Z \le z)
\]

Substitute

\[
F_Z(z) = P(X + Y \le z)
\]

\subsection{Step 2 --- Express as a Region}

The event $X+Y\le z$ defines a region in the $(x,y)$ plane.

\[
F_Z(z) =
\int\int_{x+y \le z}
f_{X,Y}(x,y)\,dx\,dy
\]

\subsection{Step 3 --- Evaluate the Integral}

The region can be written as

\[
y \le z - x
\]

Thus

\[
F_Z(z)=
\int_{-\infty}^{\infty}
\int_{-\infty}^{z-x}
f_{X,Y}(x,y)\,dy\,dx
\]

\subsection{Step 4 --- Obtain the PDF}

Differentiate the CDF

\[
f_Z(z)=\frac{d}{dz}F_Z(z)
\]

When $X$ and $Y$ are independent, this leads to the \textbf{convolution formula}

\[
f_Z(z) =
\int_{-\infty}^{\infty}
f_X(x)f_Y(z-x)\,dx
\]

\subsection{Assumption: Independent Random Variables}

If $X$ and $Y$ are independent,

\[
f_{X,Y}(x,y)=f_X(x)f_Y(y)
\]

Substituting this into the previous expression yields the convolution formula.

\section{Example Scenario Mentioned in Slides}

Let

\[
X \sim \text{Exponential}(\lambda)
\]

\[
Y \sim \text{Exponential}(\lambda)
\]

Assume independence and define

\[
Z = X + Y
\]

Using convolution,

\[
f_Z(z) =
\int_0^z f_X(x)f_Y(z-x)\,dx
\]

The exponential PDFs are

\[
f_X(x) = \lambda e^{-\lambda x}
\]

\[
f_Y(y) = \lambda e^{-\lambda y}
\]

Substituting gives

\[
f_Z(z)=
\int_0^z
\lambda e^{-\lambda x}
\lambda e^{-\lambda (z-x)}
dx
\]

Evaluating this integral produces the resulting distribution for the sum.

\section{Key Results for Exams}

\subsection{Transformation of One Random Variable}

\begin{enumerate}
\item Define the CDF
\[
F_Y(y)=P(Y\le y)
\]

\item Substitute the transformation $Y=g(X)$

\item Express the probability in terms of $X$

\item Differentiate to obtain the PDF
\end{enumerate}

\subsection{Monotonic Transformation Formula}

\[
f_Y(y) =
f_X(g^{-1}(y))
\left|
\frac{d}{dy}g^{-1}(y)
\right|
\]

\subsection{Sum of Two Random Variables}

For independent $X$ and $Y$,

\[
Z=X+Y
\]

\[
f_Z(z)=
\int_{-\infty}^{\infty}
f_X(x)f_Y(z-x)\,dx
\]

This operation is known as \textbf{convolution}.

\section{Conceptual Summary}

This lecture introduced three fundamental ideas:

\begin{itemize}
\item Transformation of random variables using the CDF method
\item Handling functions of two random variables using joint distributions
\item Computing the distribution of the sum of two variables leading to the convolution formula
\end{itemize}

These techniques are essential for deriving distributions of new variables constructed from existing random variables.

\end{document}